\section{Deep Representation Learning}\label{sec:bg:dl}

Many \gls{dl} pipelines replace hand-crafted features with a parametric map trained from data.
The central object is an encoder \(\Enc \colon \setX \to \R^{d}\) that yields an embedding \(\vect{z} = \Enc(\vect{x})\).
Later steps use \(\vect{z}\) rather than the raw coordinates.

When labels exist, \(\Enc\) can be trained so that \(\vect{z}\) separates classes.
When the target is a partition, the same encoder can be trained so that same-group points lie close and different-group points lie far.
Contrastive objectives do this by treating some pairs as positives and others as negatives.
Let \(\vect{z}_i\) and \(\vect{z}_{i}^{+}\) be normalized embeddings of a positive pair, let \(\mathrm{sim}(\cdot,\cdot)\) be cosine similarity, and let \(\tau>0\) be a temperature.
An InfoNCE-style loss is
\begin{equation}
\label{eq:bg:infonce}
\Loss_{\mathrm{NCE}}(i)
=
-\log
\frac{\exp\bigl(\mathrm{sim}(\vect{z}_i,\vect{z}_{i}^{+})/\tau\bigr)}
{\exp\bigl(\mathrm{sim}(\vect{z}_i,\vect{z}_{i}^{+})/\tau\bigr)
 + \sum_{j \neq i} \exp\bigl(\mathrm{sim}(\vect{z}_i,\vect{z}_{j}^{+})/\tau\bigr)},
\end{equation}
which is the form used, with a larger positive set, in the supervised contrastive kernel of \cref{sec:method:loss} \parencite{oord2018representation,chen2020simclr,khosla2020supcon}.

Sequence encoders based on self-attention instantiate \(\Enc\) when the input is an ordered set of tokens, here a window of \glspl{pdw}.
\Cref{sec:bg:transformer} reviews that encoder.
\Cref{sec:related:deep_clustering} surveys deep clustering methods that couple representation learning with a discrete partition.
